% !TEX program = pdflatex
% SafeBoxForest  T-RO Journal Extension
% Compile: pdflatex -> bibtex -> pdflatex x 2
\documentclass[journal]{IEEEtran}

% --- Packages ---
\usepackage[utf8]{inputenc}
\usepackage[T1]{fontenc}
\usepackage{amsmath,amssymb,amsthm}
\usepackage{graphicx}
\usepackage{algorithm}
\usepackage[noend]{algpseudocode}
\def\Or{\textbf{or}}
\newcommand{\alAnd}{\textbf{and}}
\usepackage{booktabs}
\usepackage{multirow}
\usepackage{cite}
\usepackage{xcolor}
\usepackage{hyperref}
\usepackage{cleveref}
\usepackage{bm}
\usepackage{tabularx}
\usepackage{xspace}
\usepackage{microtype}

% --- Theorem environments ---
\newtheorem{theorem}{Theorem}
\newtheorem{definition}{Definition}
\newtheorem{remark}{Remark}
\newtheorem{proposition}{Proposition}
\newtheorem{assumption}{Assumption}

% --- Custom commands ---
\newcommand{\Cfree}{\mathcal{C}_\text{free}}
\newcommand{\Cspace}{\mathcal{C}}
\newcommand{\AABB}{\text{AABB}}
\newcommand{\FK}{\text{FK}}
\newcommand{\SBF}{\text{SBF}}
\newcommand{\Real}{\mathbb{R}}
\newcommand{\calO}{\mathcal{O}}
\newcommand{\calA}{\mathcal{A}}
\newcommand{\calF}{\mathcal{F}}
\newcommand{\calR}{\mathcal{R}}
\newcommand{\bT}{\mathbf{T}}
\newcommand{\TODO}[1]{\textit{[TODO: #1]}}
\newcommand{\ie}{\textit{i.e.},\xspace}
\newcommand{\eg}{\textit{e.g.},\xspace}
\newcommand{\cf}{\textit{cf.}\@\xspace}
\newcommand{\etal}{\textit{et~al.}\@\xspace}
\DeclareMathOperator{\interior}{int}

% ==============================================================================
\begin{document}

\title{SafeBoxForest: Rapid C-space Box Cover with\\
Lifelong Envelope Caching for GCS Planning}

\author{
  TIAN,~Yu and Ren,~Hongliang
  \thanks{Department of Electronic Engineering,
  The Chinese University of Hong Kong, Hong Kong SAR, China.
  \{ty1997, hlren\}@ee.cuhk.edu.hk}
}

\markboth{IEEE Transactions on Robotics, Vol.~X, No.~Y, 2026}{}

\maketitle

% ==============================================================================
% ABSTRACT
% ==============================================================================
\begin{abstract}
We present SafeBoxForest~(SBF), a system that rapidly
covers the collision-free configuration space ($\Cfree$) of
articulated robots with certified convex boxes
for Graph-of-Convex-Sets~(GCS) motion planning.
SBF introduces a modular two-layer pipeline:
four interchangeable \emph{endpoint-AABB sources}
(IFK, CritSample, AnalyticalCrit, GCPC)
feed into three \emph{envelope representations}
(SubAABB, SubAABB\_Grid, Hull16\_Grid).
All sources produce a unified \emph{endpoint-AABB} intermediate
representation---per-link axis-aligned bounding boxes at each
interpolation point along the link---which decouples
source computation from envelope rasterisation.
Key components include:
(i)~a \emph{five-phase analytical critical-point solver}
with interval-FK certification (completeness for face-$k$, $k \leq 2$);
(ii)~a \emph{Geometric Critical-Point Cache}~(GCPC) that
pre-tabulates FK-evaluated critical configurations and
answers per-box queries $2{\times}$ faster than the analytical solver;
(iii)~the \emph{Sparse Voxel BitMask}---an $8^3$ BitBrick
tile map with turbo scanline rasterisation---producing
hull-16 envelopes $35\%$ tighter than axis-aligned bounding
boxes with constant-time collision queries; and
(iv)~the \emph{Lifelong Envelope Cache Tree}~(LECT)
with scene pre-rasterisation ($2272{\times}$ collision speedup)
and zero-loss voxel merging.
We benchmark all $4{\times}3{=}12$ pipeline configurations on a
7-DOF manipulator.
The recommended configuration (CritSample+Hull16\_Grid)
achieves $0.167$\,m$^3$ envelope volume at $0.087$\,ms
cold-start cost.
The AnalyticalCrit and GCPC sources validate CritSample
accuracy: their Hull16\_Grid volumes ($0.172$\,m$^3$) differ
by only~$3\%$, confirming that boundary sampling captures
virtually all extrema.
\end{abstract}

\begin{IEEEkeywords}
Motion planning, GCS, convex decomposition, interval arithmetic,
collision-free regions
\end{IEEEkeywords}

% ==============================================================================
\section{Introduction}\label{sec:intro}
% ==============================================================================

Motion planning for manipulators in cluttered environments
remains a fundamental challenge.
Sampling-based planners (RRT, PRM) offer general-purpose
solutions but provide no certificates and scale poorly in
narrow passages~\cite{lavalle2006planning}.
Recent work on GCS-based planning~\cite{marcucci2024motion}
formulates trajectories as convex programs over
safe regions, guaranteeing optimality and smoothness---provided
the regions accurately cover $\Cfree$.

The bottleneck shifts to \emph{region generation}.
IRIS-NP~\cite{werner2024fast} inflates ellipsoidal Regions
via semidefinite programming but requires $0.3$--$1$\,s per
region and does not reuse geometry across queries.
Clique covers~\cite{amice2024clique} offer better coverage but
are limited to polyhedral obstacle representations.

We propose SafeBoxForest~(SBF), which decomposes $\Cfree$
into axis-aligned \emph{certified} boxes.
Each box's collision-free status is verified through
interval forward kinematics (IFK) and workspace envelope
enclosures.
Key to the approach is a hierarchy of envelope
representations---from cheap axis-aligned bounding boxes to
tight sparse-voxel hull-16 grids---stored in a persistent
Lifelong Envelope Cache Tree~(LECT) that amortises computation
across boxes.

\textbf{Contributions.}
\begin{itemize}
  \item A five-phase analytical critical-point solver that computes
    provably tight envelopes with face-$k$ completeness
    for $k \leq 2$ (\cref{sec:critical}).
  \item The Sparse Voxel BitMask representation with turbo
    scanline rasterisation and constant-time collision/merge
    operations (\cref{sec:voxel}).
  \item The LECT data structure with scene pre-rasterisation
    ($2272{\times}$ collision speedup) and hull-based
    zero-loss coarsening (\cref{sec:sbf}).
  \item A Geometric Critical-Point Cache (GCPC)
    that pre-computes critical configurations for
    $2{\times}$ faster per-box queries than the analytical solver,
    validating CritSample accuracy (\cref{sec:gcpc}).
  \item A comprehensive $4{\times}3$ pipeline benchmark
    demonstrating the effectiveness of each component on
    a 7-DOF manipulator (\cref{sec:experiments}).
\end{itemize}

% ==============================================================================
\section{Related Work}\label{sec:related}
% ==============================================================================

\paragraph{Convex decomposition for motion planning.}
VCD~\cite{werner2024vcd} decomposes task space into convex polytopes
via vertex enumeration, but the combinatorial complexity grows
exponentially in dimension.
IRIS~\cite{deits2015computing} and its recent extension
IRIS-NP~\cite{werner2024fast} iteratively inflate ellipsoids
in $\Cfree$ using SDP or nonlinear programming;
each region costs $0.3$--$1$\,s and is not reusable across scenes.
Clique covers~\cite{amice2024clique} partition $\Cfree$ using
graph-theoretic methods, achieving good coverage but requiring
explicit facet representations.

\paragraph{GCS motion planning.}
Marcucci et al.~\cite{marcucci2024motion} formulates motion
planning as a shortest-path problem over a graph where
vertices are convex safe regions and edges encode transitions.
The resulting convex programs guarantee global optimality
(given the regions) and produce smooth trajectories.
SBF provides the convex regions required by GCS.

\paragraph{Interval arithmetic in robotics.}
Interval analysis has been used for verified collision
checking~\cite{merlet2004solving}, workspace
computation~\cite{merlet2006interval}, and motion
planning~\cite{jaulin2001applied}.
Our work exploits interval FK to produce \emph{certified}
bounding volumes without point-sampling, extending the
approach to multi-layer envelope representations.

\paragraph{Spatial hashing and voxel methods.}
Spatial hashing dates to Teschner et al.~\cite{teschner2003optimized}
for deformable body collision.
Recent work uses hierarchical voxel grids (VDB~\cite{museth2013vdb})
for robotic occupancy mapping.
Our BitBrick layout is lighter-weight: a single-level
$8^3$ tile map that trades hierarchy for cache-friendly
bitwise collision and zero-loss Boolean merge.

% ==============================================================================
\section{Certified Envelope Computation}\label{sec:envelope}
% ==============================================================================

% Code: include/sbf/robot/interval_fk.h, include/sbf/envelope/envelope_derive.h,
%       include/sbf/envelope/envelope_derive_critical.h

This section presents the mathematical foundations for computing
provably conservative workspace envelopes from C-space boxes.

\subsection{Interval Forward Kinematics}\label{sec:ifk}
% Code: compute_fk_full() in robot/interval_fk.h
% Code: I_sin(), I_cos(), build_dh_joint(), imat_mul_dh() in robot/interval_math.h
% Code: extract_link_aabbs() in robot/interval_fk.h

Consider a $D$-DOF serial manipulator with $L$ collision-relevant links.
Given a C-space box $B \subset \Real^D$ and obstacle AABBs $\calO$,
the goal is to certify $B \subseteq \Cfree$.

\emph{Link-AABB envelope.}
The \emph{link-AABB envelope} $\AABB_\ell(B)$ satisfies
\begin{equation}\label{eq:envelope}
  \AABB_\ell(B) \supseteq \bigcup_{\bm{q} \in B}
    \mathrm{geom}_\ell\bigl(\FK(\bm{q})\bigr).
\end{equation}
Each link~$\ell$ is modeled as a capsule: a line segment between joint
origins $\bm{p}_\ell^{\mathrm{prox}}$ and $\bm{p}_\ell^{\mathrm{dist}}$,
inflated by radius~$r_\ell$.

\begin{assumption}[Geometry Enclosure]\label{asm:geom}
For each link~$\ell$, the capsule of radius~$r_\ell$ encloses the
true mesh geometry at every configuration.
\end{assumption}

Since every interior point of the capsule segment is a convex combination
of its two endpoints, coordinate extremes over all configurations in~$B$
are attained at the endpoints alone.
The AABB therefore reduces to
\begin{equation}\label{eq:endpoint_aabb}
  \AABB_\ell(B)
  = \mathrm{bbox}\!\Bigl(
      \bigl[\bm{p}_\ell^{\mathrm{prox}}\bigr]_B,\;
      \bigl[\bm{p}_\ell^{\mathrm{dist}}\bigr]_B
    \Bigr) \oplus r_\ell,
\end{equation}
where $[\cdot]_B$ is the interval enclosure over $B$ and $\oplus\,r_\ell$
is a uniform inflation.

\emph{Computing via interval-FK.}
Replacing scalar joint angles with interval counterparts in the DH chain
and multiplying sequentially,
\begin{equation}\label{eq:chain}
  [\bT^{0:k}] = [\bT^{0:k-1}] \otimes [\bT_k],
\end{equation}
yields interval transforms for all endpoints.
The two endpoint positions in~\eqref{eq:endpoint_aabb} are read directly
from the translation components of $[\bT^{0:\ell-1}]$ and $[\bT^{0:\ell}]$;
no additional FK evaluations are needed.
By inclusion-monotonicity of interval
arithmetic~\cite{moore2009introduction}, the result is a
provably conservative enclosure; chain multiplication may over-approximate
(the \emph{wrapping effect}), affecting tightness but not soundness.
By construction, over-approximation can only cause \emph{false rejections}
(a free box mistakenly rejected), never false acceptances, so every
box that passes the envelope check is guaranteed collision-free.

\emph{Active-link filtering.}
Not all links contribute to collision geometry.
Zero-length links (coincident joint origins) produce degenerate AABBs
and are automatically excluded from the envelope and collision pipeline.
Only $L_\text{act} \leq L$ \emph{active links} are retained, reducing
both storage and collision-check cost.
The active-link map is determined once from the robot model and stored
as a compact index array.

\subsection{Subdivision Enhancements}\label{sec:subdivision}
% Code: derive_aabb_subdivided() in envelope/envelope_derive.h
% Code: compute_fk_incremental() for prefix reuse

Interval-FK can over-approximate due to the wrapping effect.
Two complementary subdivision strategies reduce this conservatism while
preserving certification.

\emph{C-space box subdivision.}
Bisecting box~$B$ along a dimension into $B_1$ and $B_2$ and computing
their envelopes separately yields tighter enclosures than treating $B$
as a whole.
Recursive $k$-level subdivision produces $2^k$ sub-boxes.
Certification is preserved by inclusion-monotonicity:
each sub-box envelope is a valid over-approximation, and their union
is contained in the original ($\bigcup_i \AABB_\ell(B_i) \subseteq \AABB_\ell(B)$)
but pointwise tighter.
Sub-boxes share FK prefix chains up to the split dimension,
so the extra cost is sub-linear.

\emph{Incremental FK.}
When a parent node splits along dimension~$d$, the interval transform
prefix $[\bT^{0:d-1}]$ is identical for both children.
Only the suffix $[\bT^{d:D}]$ must be recomputed.
On a 7-DOF manipulator, this saves approximately 43\% of FK cost on
average; the saving grows with joint index~$d$.

\emph{Link subdivision.}
For a long link whose AABB spans a large workspace volume,
splitting link~$\ell$ into $n$ equal segments and computing per-segment
AABBs independently yields a tighter aggregate envelope.
The $n{-}1$ intermediate segment endpoints are obtained by linearly
interpolating the proximal and distal joint-origin positions
$\bm{p}_\ell^{\mathrm{prox}}$ and $\bm{p}_\ell^{\mathrm{dist}}$,
which are already available from the FK chain; no additional FK
evaluation is required.
Certification is preserved: each sub-segment AABB encloses the
corresponding mesh portion under \cref{asm:geom}.

\subsection{Endpoint-AABB: Unified Intermediate Representation}%
\label{sec:endpoint_aabb}

All envelope derivation methods produce a common intermediate
representation called the \emph{endpoint-AABB}.
For each active link~$\ell$ and subdivision resolution $n$,
the endpoint-AABB is the array of $n{+}1$ geometry-only
(un-inflated) axis-aligned bounding boxes
\begin{equation}\label{eq:ep_aabb}
  E_{\ell,s} \;=\;
  \mathrm{bbox}\Bigl(
    \bigl[\bm{p}_\ell^\mathrm{prox}(s/n)\bigr]_B,\;
    \bigl[\bm{p}_\ell^\mathrm{dist}(s/n)\bigr]_B
  \Bigr),
  \qquad s = 0,1,\ldots,n,
\end{equation}
where $[\cdot]_B$ denotes the interval enclosure over the C-space
box~$B$ and $(s/n)$ is a linear interpolation parameter along the
link segment.
The total storage is $L_\text{act} \times (n{+}1) \times 6$ floats.

The key property is that the envelope representation (SubAABB,
SubAABB\_Grid, or Hull16\_Grid) is derived \emph{solely} from
the endpoint-AABB array---it does not depend on how the
endpoint-AABBs were produced.
This decouples the \emph{source} computation
(IFK, CritSample, AnalyticalCrit, or GCPC) from the
\emph{envelope rasterisation}:
\begin{equation}\label{eq:pipeline}
  \underbrace{\text{Source}}_{\text{IFK / CritSample / Analytical / GCPC}}
  \;\xrightarrow{\;E_{\ell,s}\;}
  \underbrace{\text{Envelope}}_{\text{SubAABB / SubAABB\_Grid / Hull16\_Grid}}.
\end{equation}

\emph{Deriving sub-AABBs from endpoint-AABBs.}
Given target subdivision count $n_\text{tgt}$ (which must divide $n$),
the $k$-th sub-AABB for link~$\ell$ is
\begin{equation}
  A_{\ell,k} = \bigcup_{s = k\cdot r}^{(k+1)\cdot r}\!E_{\ell,s}
  \;\oplus\; r_\ell,
  \qquad r = n / n_\text{tgt},
\end{equation}
recovering the subdivided link envelope of \cref{sec:subdivision}
with per-link radius inflation added at the envelope layer.

\emph{Deriving Hull16\_Grid from endpoint-AABBs.}
The proximal box $E_{\ell,0}$ and distal box $E_{\ell,n}$
define the 16-vertex convex hull of \cref{sec:hull16}.
The hull-16 turbo rasteriser takes these two endpoint AABBs
(plus link radius) directly, without passing through the
sub-AABB intermediate.

\subsection{Analytical Critical-Point Solver}\label{sec:critical}
% Code: derive_aabb_critical_analytical() in envelope/envelope_derive_critical.h
% Code: AnalyticalCriticalConfig in envelope/envelope_derive_critical.h
% Exp-28: experiments/28_gcpc_validation.cpp

Interval-FK provides a \emph{certified} envelope but may over-approximate
due to the wrapping effect.
We complement it with a \emph{tight} envelope obtained by solving for
the true coordinate extremes analytically.
We first establish structural properties of DH kinematics that
enable exact solutions, then present the 5-phase solver and prove
its completeness guarantees.

% ---------- structural properties ----------

\subsubsection{Trigonometric Structure of DH Chains}\label{sec:critical_structure}

Each $4{\times}4$ homogeneous DH matrix $\bT_j$ for a revolute joint~$j$
is affine in $(\cos q_j,\,\sin q_j)$:
\begin{equation}\label{eq:dh_affine}
  \bT_j(q_j)
  = \bm{M}_j^{(c)}\cos q_j
  + \bm{M}_j^{(s)}\sin q_j
  + \bm{M}_j^{(0)},
\end{equation}
where $\bm{M}_j^{(c)}, \bm{M}_j^{(s)}, \bm{M}_j^{(0)}$ are constant
matrices determined by the DH parameters $(a_j,\alpha_j,d_j)$.
The composed transform is
$\bT^{0:\ell} = \bT_1 \cdots \bT_\ell$
and the position coordinate of the distal origin is
$p_{\ell,d}(\bm{q}) = [\bT^{0:\ell}]_{d,4}$.
Since each factor is affine in its own $(c_j,s_j)$ pair,
the product is \emph{multilinear} in
$\{(\cos q_j,\sin q_j)\}_{j=1}^\ell$.

\begin{proposition}[Univariate sinusoidal]\label{prop:sinusoid}
Fix all joints except~$j$.
Then
\begin{equation}\label{eq:sin_form}
  p_{\ell,d}(q_j)
  = A_j\cos q_j + B_j\sin q_j + C_j,
\end{equation}
where $A_j, B_j, C_j$ depend only on the fixed joints and
DH parameters.
This function has \emph{exactly two} critical points per $2\pi$
period: $q_j^* = \mathrm{atan2}(B_j,A_j)$ (global max)
and $q_j^* + \pi$ (global min).
\end{proposition}
\begin{proof}
By~\eqref{eq:dh_affine}, $\bT_j$ is the only factor containing
$q_j$; all other factors are constant matrices.
Collecting terms in the matrix product yields~\eqref{eq:sin_form}.
The derivative
$\partial p/\partial q_j = -A_j\sin q_j + B_j\cos q_j = 0$
gives $\tan q_j = B_j/A_j$, solved uniquely (mod $\pi$) by $\mathrm{atan2}$.
\end{proof}

\begin{proposition}[Bivariate bilinear trigonometric]\label{prop:bilinear}
Fix all joints except $(i,j)$.
Then
\begin{equation}\label{eq:bilinear}
  p_{\ell,d}(q_i,q_j) = \sum_{m=1}^{9} a_m\,\phi_m(q_i,q_j),
\end{equation}
where $\{\phi_m\} = \{c_i c_j,\; c_i s_j,\; s_i c_j,\; s_i s_j,\;
c_i,\; s_i,\; c_j,\; s_j,\; 1\}$ and the nine coefficients $a_m$
depend on the fixed joints.
\end{proposition}
\begin{proof}
$\bT_i$ and $\bT_j$ are distinct factors in the matrix product
$\bT^{0:\ell}$.
Each contributes $(c_k, s_k)$ linearly; their product introduces
exactly the four cross-terms $c_i c_j, c_i s_j, s_i c_j, s_i s_j$
plus the six lower-order and constant terms.
\end{proof}

\begin{proposition}[Bivariate critical-point resultant]\label{prop:resultant}
Setting $\nabla_{(q_i,q_j)} p_{\ell,d} = \bm{0}$ in the bilinear
model~\eqref{eq:bilinear} and applying the Weierstrass substitution
$t_j = \tan(q_j/2)$ reduces the system to a single univariate
polynomial of degree~$\leq 8$ in~$t_j$.
All real roots can therefore be found exactly via the companion-matrix
eigenvalue method.
\end{proposition}
\begin{proof}
Setting $\partial p / \partial q_i = 0$ yields
$\tan q_i = P(q_j) / Q(q_j)$ where $P,Q$ are of the
form $\alpha c_j + \beta s_j + \gamma$.
Substituting into $\partial p / \partial q_j = 0$ and clearing
denominators via $c_j = (1-t_j^2)/(1+t_j^2)$,
$s_j = 2t_j/(1+t_j^2)$ gives a rational equation whose numerator
is a polynomial in $t_j$ of degree~$\leq 8$.
The companion matrix of this polynomial has eigenvalues at all roots,
including complex ones; only real eigenvalues correspond to
physical configurations.
\end{proof}

% ---------- 5-phase solver ----------

\subsubsection{Five-Phase Solver}\label{sec:five_phase}

We now describe the solver.
Consider maximising (or minimising) $p_{\ell,d}(\bm{q})$ over the box
$B = \prod_{j=1}^{D} [q_j^-,\, q_j^+]$.
By compactness, the extremum exists.
By the Karush--Kuhn--Tucker (KKT) conditions on the box constraints,
at any extremiser~$\bm{q}^*$ the joints partition into
a \emph{free set}
$F = \{j : q_j^* \in (q_j^-,q_j^+)\}$ and a
\emph{boundary set} $\bar F = \{j : q_j^* \in \{q_j^-,q_j^+\}\}$,
with $\partial p_{\ell,d}/\partial q_j(\bm{q}^*) = 0$ for every
$j \in F$.
In other words, the extremum lies on a \emph{face} of~$B$ of
dimension~$|F|$, and satisfies the gradient-zero condition on
that face.
The five phases enumerate candidate extremisers by face dimension:

\emph{Phase~0 -- Vertex enumeration ($|F|=0$).}
The candidate set for each joint is
$\mathcal{K}_j = \{q_j^-,\, q_j^+\}
\cup \{ k\pi/2 : k\in \mathbb{Z},\;
q_j^- < k\pi/2 < q_j^+ \}$,
and all configurations in the Cartesian product
$\mathcal{K}_1 \times \cdots \times \mathcal{K}_D$ are evaluated.
This set \emph{contains} all $2^D$ true vertices of~$B$
(each joint at its endpoint), plus interior $k\pi/2$ values that
serve as high-quality background seeds for later phases.

\emph{Phase~1 -- Edge critical points ($|F|=1$).}
For each joint~$j$ in turn and each of the~$2^{D-1}$ boundary
assignments of $\bar F = \{1,\ldots,D\}\setminus\{j\}$,
the 3-point sinusoidal fit of \cref{prop:sinusoid} determines
$A_j, B_j, C_j$ exactly, and the two critical points
$q_j^* = \mathrm{atan2}(B_j,A_j)$ and $q_j^*+\pi$ are tested.
Since~\eqref{eq:sin_form} has \emph{no other} critical points within
any one period, this procedure is exhaustive over all 1-faces.

\emph{Phase~2 -- Face critical points ($|F|=2$).}
For each joint pair $(i,j)$ and each boundary assignment of the
remaining $D{-}2$ joints, the 9-coefficient bilinear trig model
(\cref{prop:bilinear}) is fitted exactly from a $3{\times}3$ grid
sample.
The degree-$\leq 8$ resultant (\cref{prop:resultant}) is then
solved via the companion matrix; every real root yields
a candidate $(q_i^*,q_j^*)$ pair.
By construction, all critical points of~\eqref{eq:bilinear}
inside~$B$ are recovered.

\emph{Phase~2.5 -- Pair-constrained critical points.}
\begin{proposition}[Pair-manifold reduction]\label{prop:pair_reduce}
On the affine manifold $q_i + q_j = C$ with $C = k\pi/2$,
substituting $q_j = C - q_i$ into~\eqref{eq:bilinear} and fixing
all other joints reduces $p_{\ell,d}$ to the sinusoidal
form~\eqref{eq:sin_form} in~$q_i$ alone.
\end{proposition}
\begin{proof}
$\cos(C-q_i) = \cos C \cos q_i + \sin C \sin q_i$ and
$\sin(C-q_i) = \sin C \cos q_i - \cos C \sin q_i$ are linear
in $(\cos q_i, \sin q_i)$.
Substituting into every $(c_j, s_j)$ term of~\eqref{eq:bilinear}
collapses the bilinear form to $A'\cos q_i + B'\sin q_i + C'$.
\end{proof}
Phase~2.5a exploits this by enumerating all integer~$k$ with
$q_i^- + q_j^- \leq k\pi/2 \leq q_i^+ + q_j^+$,
substituting $q_j = k\pi/2 - q_i$,
and solving the resulting sinusoidal exactly as in Phase~1.
This captures critical points on manifolds invisible to the
per-joint Phase~1 solver.

\emph{Phase~3 -- Interior coordinate descent ($|F|\geq 3$).}
For faces of dimension~$\geq 3$ the KKT system is a system of
$|F|$ coupled trigonometric equations.
We solve it approximately by block-coordinate descent:
in each sweep, every free joint~$j$ is optimised exactly
(via \cref{prop:sinusoid}) with all others held fixed, cycling
repeatedly.
Multi-start initialisation from Phase~0/1/2/2.5 seeds and
pair-coupled updates ($q_j \leftarrow k\pi/2 - q_i$ after
optimising~$q_i$) improve convergence.

% ---------- completeness ----------

\subsubsection{Completeness Guarantees}\label{sec:completeness}

\begin{theorem}[Face-$k$ completeness for $k \leq 2$]\label{thm:face_complete}
Let $\bm{q}^*$ be a global extremiser of $p_{\ell,d}$ over
$B = \prod_j [q_j^-,q_j^+]$ with free set $F$ satisfying $|F|\leq 2$.
Then the analytical solver finds a configuration with equal or
better objective value, provided that the boundary enumeration
covers all $2^{|\bar F|}$ assignments of $\bar F$.
\end{theorem}
\begin{proof}
$|F|=0$:
$\bm{q}^*$ has every joint at a boundary value.
Since $\mathcal{K}_j \supseteq \{q_j^-,q_j^+\}$, the Cartesian product
in Phase~0 includes~$\bm{q}^*$.

$|F|=1$:
The single free joint~$j$ satisfies the sinusoidal
equation $\partial p/\partial q_j = 0$ with all other joints at
boundary values.
Phase~1 enumerates all $2^{D-1}$ boundary assignments and, for each,
finds \emph{all} critical points of the sinusoidal
by \cref{prop:sinusoid}.
Hence $\bm{q}^*$ is found.

$|F|=2$:
The two free joints $(i,j)$ satisfy
$\nabla_{(q_i,q_j)} p_{\ell,d} = \bm{0}$ with the
remaining joints at boundary values.
Phase~2 enumerates all $2^{D-2}$ boundary assignments and, for each,
finds all critical points of the bilinear-trig function via
\cref{prop:resultant}.
Hence $\bm{q}^*$ is found.
\end{proof}

\begin{remark}[Boundary budget]\label{rem:budget}
\cref{thm:face_complete} requires exhaustive boundary enumeration.
In our implementation, the budget is $2^{D-1} \leq 128$ for Phase~1
and $2^{D-2} \leq 64$ for Phase~2.
For a 7-DOF manipulator ($D=7$), this gives $2^6=64$ and $2^5=32$,
well within budget; the theorem therefore holds unconditionally for
$D\leq 8$.
\end{remark}

\begin{theorem}[Pair-manifold completeness]\label{thm:pair_complete}
Let $\bm{q}^*$ be a global extremiser with
$|F| \leq 2$ and $q_i^* + q_j^* = k\pi/2$ for some integer~$k$ and
some pair~$(i,j)$.
Then Phase~2.5a finds a configuration with equal or better objective
value.
\end{theorem}
\begin{proof}
By \cref{prop:pair_reduce}, substituting the constraint reduces the
problem to a sinusoidal in~$q_i$.
Phase~2.5a enumerates all valid $k$ and, for each, solves the
sinusoidal exactly.
\end{proof}

% ---------- dual phase 3 ----------

\subsubsection{Fused Dual-Phase-3 Mode}\label{sec:dual_phase3}

Phase~3 is inherently heuristic: block-coordinate descent on a
non-convex multivariate trigonometric function can converge to a
sub-optimal local extremum.
We identify a systematic failure mode and eliminate it.

\emph{Basin-shift problem.}
Phase~2's companion-matrix solutions update the running best
configuration~$\bm{q}^\text{best}$, which Phase~3 uses as its
starting point.
This update may move $\bm{q}^\text{best}$ into the basin of
attraction of a \emph{worse} local extremum for the coordinate-descent
process---a phenomenon we call \emph{basin shift}.
Removing Phase~2 entirely would prevent the shift but sacrifices the
2-face critical points that Phase~2 provides.

\emph{Solution: dual-phase-3 with checkpoint.}
We checkpoint the per-link extremes after Phase~1
($E^{01}$, before Phase~2 modifies the state) and proceed normally
through Phases~2--3 to obtain $E^{0123}$.
A second Phase~3 pass then starts from the checkpoint $E^{01}$,
yielding $E^{013}$.
The final result is the element-wise merge:
\begin{equation}\label{eq:dual_merge}
  E^*_{\ell,d} = \bigl[\,
    \min\!\bigl(E^{0123}_{\ell,d,\mathrm{lo}},\;
                E^{013}_{\ell,d,\mathrm{lo}}\bigr),\;
    \max\!\bigl(E^{0123}_{\ell,d,\mathrm{hi}},\;
                E^{013}_{\ell,d,\mathrm{hi}}\bigr)
  \,\bigr].
\end{equation}

\begin{proposition}[Merge dominance]\label{prop:merge}
$E^*$ is at least as tight as both $E^{0123}$ and $E^{013}$:
for every link~$\ell$, dimension~$d$,
$E^*_{\ell,d,\mathrm{lo}} \leq \min(E^{0123}_{\ell,d,\mathrm{lo}},\,
E^{013}_{\ell,d,\mathrm{lo}})$
and
$E^*_{\ell,d,\mathrm{hi}} \geq \max(E^{0123}_{\ell,d,\mathrm{hi}},\,
E^{013}_{\ell,d,\mathrm{hi}})$.
If a global extremum is found by \emph{either} Phase~3 run,
it appears in~$E^*$.
\end{proposition}
\begin{proof}
Follows directly from the $\min/\max$ definition of~\eqref{eq:dual_merge}.
\end{proof}

\begin{remark}[Overhead]
Checkpointing $E^{01}$ requires one deep copy of the per-link
extremes array (negligible).
The second Phase~3 pass adds one coordinate-descent run but
\emph{avoids} repeating Phases~0 and~1.
Compared with the na\"ive approach of running the full Phases~0--3
pipeline twice (once with Phase~2 and once without), the fused
mode saves two redundant Phase~0+1 computations, yielding a
$15.3\times$ speedup (43\,ms vs.\ 660\,ms; Exp-28, Panda + IIWA14).
\end{remark}

% ---------- interval-FK certification ----------

\subsubsection{Interval-FK Certification}\label{sec:ifk_cert}

Phases~0--2.5 are provably complete for $|F|\leq 2$
(\cref{thm:face_complete,thm:pair_complete}).
For $|F|\geq 3$, Phase~3 is a heuristic that may miss the global
extremum.
Dual-phase-3 mitigates this empirically (0/180 misses), but does
not constitute a proof.
We close this theoretical gap via \emph{interval-FK certification}:
a lightweight post-processing step that restores soundness
unconditionally.

\emph{Observation.}
The analytical solver evaluates FK at specific configurations; its
output AABB satisfies
$\AABB^{\mathrm{anal}}_\ell \subseteq \AABB^{\mathrm{true}}_\ell$
(a lower bound on the true extent), with equality when all critical
points are found.
Interval-FK (\cref{sec:ifk}) produces a guaranteed \emph{outer}
bound
$\AABB^{\mathrm{ifk}}_\ell \supseteq \AABB^{\mathrm{true}}_\ell$.
The true AABB is sandwiched:
\begin{equation}\label{eq:sandwich}
  \AABB^{\mathrm{anal}}_\ell
  \;\subseteq\; \AABB^{\mathrm{true}}_\ell
  \;\subseteq\; \AABB^{\mathrm{ifk}}_\ell.
\end{equation}

\emph{Certification step.}
After the analytical pipeline writes its output AABB, a single
interval-FK evaluation computes the outer bound per link.
For each AABB face, the output is clamped to be at least as
extreme as the interval-FK bound:
\begin{equation}\label{eq:cert_clamp}
  \widehat{\AABB}_{\ell,d}^{\mathrm{lo}}
    = \min\!\bigl(\AABB_{\ell,d}^{\mathrm{anal,lo}},\;
                   \AABB_{\ell,d}^{\mathrm{ifk,lo}}\bigr),
  \quad
  \widehat{\AABB}_{\ell,d}^{\mathrm{hi}}
    = \max\!\bigl(\AABB_{\ell,d}^{\mathrm{anal,hi}},\;
                   \AABB_{\ell,d}^{\mathrm{ifk,hi}}\bigr).
\end{equation}

\begin{theorem}[Certified soundness]\label{thm:cert_sound}
With certification enabled, the output AABB satisfies
$\widehat{\AABB}_\ell \supseteq \AABB^{\mathrm{true}}_\ell$
for every active link~$\ell$ and sub-segment, regardless of
whether Phase~3 found all critical points.
\end{theorem}
\begin{proof}
By~\eqref{eq:sandwich},
$\AABB_{\ell,d}^{\mathrm{ifk,lo}} \leq
 \AABB_{\ell,d}^{\mathrm{true,lo}}$
and
$\AABB_{\ell,d}^{\mathrm{ifk,hi}} \geq
 \AABB_{\ell,d}^{\mathrm{true,hi}}$.
The $\min/\max$ clamp in~\eqref{eq:cert_clamp} therefore
ensures $\widehat{\AABB}_{\ell,d}^{\mathrm{lo}} \leq
\AABB_{\ell,d}^{\mathrm{true,lo}}$
and $\widehat{\AABB}_{\ell,d}^{\mathrm{hi}} \geq
\AABB_{\ell,d}^{\mathrm{true,hi}}$
for every $d$.
\end{proof}

\begin{remark}[Gap metric]
The per-face gap
$\delta_{\ell,d}
  = \AABB^{\mathrm{ifk}}_{\ell,d} - \AABB^{\mathrm{anal}}_{\ell,d}$
quantifies the interval-FK wrapping effect on that face.
When $\delta = 0$, both bounds agree and the result is
\emph{simultaneously sound and tight}.
In our experiments, dual-phase-3 achieves $\delta = 0$
(to float precision) on all faces for narrow SBF-typical boxes;
the residual gap appears only on wide boxes (20\% of interval range)
and is due to interval-FK conservatism, not analytical incompleteness.
\end{remark}

\begin{remark}[Overhead]
The certification step requires one interval-FK call per box ($< 5\,\mu$s)
and a per-face comparison loop.
On Panda 7-DOF with 60 random trials (80\% narrow, 20\% wide boxes),
the certified mode adds $< 0.3$\,ms average overhead on top of the
${\sim}28$\,ms analytical solver ($< 1$\% relative cost).
All 120 trials (Panda + IIWA14) pass soundness verification with zero
violations.
\end{remark}

\subsection{Geometric Critical-Point Cache (GCPC)}\label{sec:gcpc}
% Code: include/sbf/envelope/gcpc_cache.h, src/envelope/gcpc_cache.cpp
% Exp: experiments/exp_pipeline_benchmark.cpp

The analytical solver of \cref{sec:critical} is provably complete
for face-$k$ critical points ($k \leq 2$), but its per-box cost
($\sim$\,$150$\,ms, IIWA14) is dominated by the combinatorial
boundary enumeration.
We introduce a \emph{Geometric Critical-Point Cache}~(GCPC) that
amortises this cost across all queries for a given robot.

\emph{One-time precomputation.}
For each active link~$\ell$, every $k\pi/2$ critical configuration
(and its composite-angle variants) is enumerated within the full
joint limits.
For each configuration, the FK position is evaluated and stored as
a \texttt{GcpcPoint} with fields
$(\ell,\, q_\text{eff},\, A,\, B,\, C,\, R)$;
points are organised into a per-link KD-tree in the joint subspace.
An enrichment pass adds interior critical points discovered
via coordinate descent from random seeds.
For IIWA14 (7-DOF), the cache contains $\sim$\,$16{,}000$ points
and is built in $\sim$\,$300$\,ms.

\emph{Per-box query.}
Given a C-space box~$B$, the KD-tree range query retrieves all
cached points whose effective joints fall within~$B$.
Each retrieved point contributes to the running min/max
of the per-link position coordinates.
The output is the same endpoint-AABB format
(\cref{sec:endpoint_aabb}) as the analytical solver:
$L_\text{act} \times (n{+}1) \times 6$ floats.

\emph{Cost--tightness trade-off.}
GCPC produces endpoint-AABBs at $72$\,ms per box---$2.1{\times}$
faster than the analytical solver ($155$\,ms)---while achieving
identical envelope volumes (both $0.172$\,m$^3$ for Hull16\_Grid,
\cref{tab:pipeline}).
The cache is \emph{scene-independent} and can be serialised to
disk for reuse across sessions.

% ==============================================================================
\section{Sparse Voxel BitMask}\label{sec:voxel}
% ==============================================================================

% Code: include/sbf/voxel/bit_brick.h, include/sbf/voxel/voxel_grid.h,
%       include/sbf/voxel/hull_rasteriser.h, src/voxel/hull_rasteriser.cpp
% Experiment: experiments/voxel_proto/main.cpp

The AABB envelope of \cref{sec:ifk} represents each link as a single
axis-aligned box~\eqref{eq:endpoint_aabb}.
While cheap, this representation wastes volume at joints with
large angular range: the AABB of a \emph{curved} swept arm includes
substantial dead space that inflates false-positive collision rates.
GCS-based planning requires convexity only in C-space; the workspace
occupancy representation may be arbitrarily non-convex.
We exploit this by replacing per-link AABBs with a
\emph{Sparse Voxel BitMask} that tightly fills the true
swept-capsule geometry.

\subsection{BitBrick Layout and Spatial Hashing}\label{sec:bitbrick}

The workspace is decomposed into $8^3{=}512$-voxel tiles
(\emph{BitBricks}), each stored as $8{\times}\mathtt{uint64}$
(64~bytes, one cache line).
The bit at position $(x,y,z)$ occupies word~$z$, bit~$8y{+}x$.
Tiles are addressed by an integer $(b_x,b_y,b_z)$ triple
and stored in a hash map (FNV-1a keyed), so only workspace regions
actually reached by the manipulator allocate memory.
For a typical 7-DOF IIWA14 arm at $\Delta{=}0.02$\,m resolution,
a medium-box envelope occupies $\sim$50 bricks ($\approx$3.2\,KB),
fitting entirely in L1 cache.

\subsection{16-Point Convex Hull Rasterisation}\label{sec:hull16}

Given proximal endpoint box~$B_1$ and distal endpoint box~$B_2$
for a link sub-segment, the swept volume of the axis is exactly
the convex hull of the two boxes' corner sets%
\footnote{Because every swept point $(1{-}t)\bm{p}_1+t\bm{p}_2$
lies inside $\mathrm{Conv}(\mathrm{corners}(B_1)\cup\mathrm{corners}(B_2))$.}:
\begin{equation}\label{eq:hull16}
  S_\ell^{(s)} = \mathrm{Conv}\bigl(
    \mathrm{corners}(B_1^{(s)}) \cup \mathrm{corners}(B_2^{(s)})
  \bigr),
\end{equation}
which has at most 16 vertices (8 from each box).
At parameter~$t \in [0,1]$, the cross-section
$B(t) = (1{-}t)\,B_1 + t\,B_2$ is an AABB whose bounds are
\emph{linear} in~$t$:
\begin{equation}\label{eq:bt_linear}
  \mathrm{lo}_a(t) = b_{1,\mathrm{lo}}^a + t\,\Delta\mathrm{lo}^a, \qquad
  \mathrm{hi}_a(t) = b_{1,\mathrm{hi}}^a + t\,\Delta\mathrm{hi}^a,
\end{equation}
for each axis~$a \in \{x,y,z\}$.
The squared distance from a query point~$\bm{q}$ to the hull is
$\min_{t \in [0,1]} \mathrm{dist}^2(\bm{q}, B(t))$, which is a
\emph{convex piecewise-quadratic} function of~$t$ with at most
6 breakpoints (2 per axis), solvable analytically in $O(1)$.

To certify conservation, voxels are marked whenever their
centre lies within distance $r_\ell + \sqrt{3}\,\Delta/2$ of
the hull, where $\sqrt{3}\,\Delta/2$ is the half-diagonal of a
voxel cube---this ensures every continuous point inside the hull
falls within at least one marked voxel.

\subsection{Turbo Scanline Algorithm}\label{sec:turbo}

Naively testing every voxel in the bounding box costs
$\calO(n_x\,n_y\,n_z)$.
We eliminate all per-voxel distance evaluations by exploiting the
linear structure of~\eqref{eq:bt_linear}.

\emph{YZ $t$-range solver.}
For each $(y,z)$ scanline, we solve the piecewise-quadratic inequality
\begin{equation}\label{eq:yz_ineq}
  \mathrm{dist}^2_{yz}\!\bigl((q_y,q_z),\, B(t)|_{yz}\bigr) \;\le\; r_{\text{eff}}^2
\end{equation}
to obtain the active $t$-range $[t_{\mathrm{lo}}, t_{\mathrm{hi}}]$.
The inequality is convex in~$t$, so the feasible set is a single
interval (or empty); its endpoints are obtained by solving at most
5 quadratic segments via the discriminant, in $O(1)$ per scanline.

\emph{Conservative $x$-fill.}
Since $\mathrm{lo}_x(t)$ and $\mathrm{hi}_x(t)$ are linear in~$t$,
the extremes over $[t_{\mathrm{lo}}, t_{\mathrm{hi}}]$ are attained
at the endpoints:
\begin{equation}\label{eq:turbo_x}
  x \in \bigl[\,
    \min\!\bigl(\mathrm{lo}_x(t_{\mathrm{lo}}),\,\mathrm{lo}_x(t_{\mathrm{hi}})\bigr) - r_{\text{eff}},\;
    \max\!\bigl(\mathrm{hi}_x(t_{\mathrm{lo}}),\,\mathrm{hi}_x(t_{\mathrm{hi}})\bigr) + r_{\text{eff}}
  \,\bigr].
\end{equation}

\begin{proposition}[Conservation]\label{prop:turbo_conserv}
If point $P{=}(p_x,p_y,p_z)$ lies inside the inflated hull
$S_\ell^{(s)} \oplus r_{\text{eff}}$, then there exists
$t^* \in [t_{\mathrm{lo}},t_{\mathrm{hi}}]$ with
$\mathrm{dist}(P,B(t^*)) \le r_{\text{eff}}$, and therefore
$p_x$ lies within the range~\eqref{eq:turbo_x} by linearity of
$\mathrm{lo}_x$ and $\mathrm{hi}_x$.
\end{proposition}

\emph{Batch bit fill.}
Filling the contiguous $x$-range within a single BitBrick
$(b_x, b_y, b_z)$ reduces to a single 64-bit OR with a
precomputed bit mask, instead of per-voxel set operations.
Across bricks, the fill advances in 8-voxel strides.
The overall complexity is $\calO(n_y \cdot n_z)$ with $O(1)$ per
scanline, plus $O(n_x/8)$ for the batch fill.

\subsection{Voxel Collision and Merge}\label{sec:voxel_ops}

\emph{Collision detection.}
Testing collision between a robot envelope grid $G_R$ and an obstacle
grid $G_O$ iterates over the \emph{smaller} hash map.
For each matching tile coordinate, a single 64-bit AND per word
detects overlap:
\begin{equation}
  \text{collide}(G_R, G_O) \;\Leftrightarrow\;
  \exists\, \bm{b}:\;
  \bigl(G_R[\bm{b}].\mathrm{words}[k]
  \;\mathbin{\&}\; G_O[\bm{b}].\mathrm{words}[k]\bigr) \ne 0.
\end{equation}
Early exit on the first nonzero AND makes collision testing
$O(1)$ in the expected case.

\emph{Zero-loss merge.}
Merging two grids (forest node union) is an exact per-brick OR:
$G_{\text{parent}}[\bm{b}] = G_{\text{left}}[\bm{b}] \mathbin{|} G_{\text{right}}[\bm{b}]$.
Unlike AABB union---which inflates to the bounding box of two AABBs
and introduces dead space---voxel OR preserves the exact shape union
with zero volume inflation.

\emph{Preliminary benchmarks (IIWA14, 7-DOF, medium 10\% box,
$\Delta{=}0.02$\,m).}
The unified \texttt{fill\_hull16} rasteriser derives the 16-point
convex hull directly from each link's proximal and distal
endpoint-AABBs (plus link radius), bypassing the sub-AABB
intermediate entirely.
The turbo scanline rasterises 4 active links at
$\approx\!230$\,\textmu s, within 7\% of the sub-AABB baseline
($\approx\!217$\,\textmu s).
Compared to the brute-force hull rasteriser ($\approx\!2160$\,\textmu s),
this is a $9.4{\times}$ speedup.
The hull envelope is $\sim\!27\%$ tighter than the sub-AABB
envelope for small-to-medium boxes (measured by voxel volume ratio
to AABB-sum volume: 73\% vs.\ 79\%).
Voxel-voxel collision takes $\approx\!3.3$\,ns for $\sim\!50$-brick
envelopes, and grid merge costs $\approx\!3.9$\,\textmu s.

% ==============================================================================
\section{SafeBoxForest Architecture}\label{sec:sbf}
% ==============================================================================

% Code: include/sbf/envelope/frame_store.h, include/sbf/envelope/grid_store.h
%       include/sbf/forest/, src/forest/

This section describes the system architecture that integrates
the envelope computation (\cref{sec:envelope}) and voxel representation
(\cref{sec:voxel}) into a scalable C-space box decomposition.

\subsection{Lifelong Envelope Cache Tree (LECT)}\label{sec:lect}
% Code: FrameStore class in envelope/frame_store.h (SoA layout, mmap persistence)
% Code: GridStore class in envelope/grid_store.h (bitfield grid cache)

The LECT is a lazy KD-tree over C-space that caches joint-origin
interval vectors and derives envelopes on demand,
exploiting both subdivision techniques with near-zero overhead.

\emph{Structure.}
The root covers the full configuration space $\Cspace$.
Each internal node bisects one dimension.
Each node~$v$ stores a C-space box $B_v$ and $L$ joint-origin
interval vectors $[\bm{p}_1]_{B_v},\ldots,[\bm{p}_L]_{B_v}$,
where $\bm{p}_\ell$ is the distal joint origin of link~$\ell$.
The base origin $\bm{p}_0$ is a fixed constant and is never stored.
The link envelope is derived on demand---either as an AABB:
\begin{equation}\label{eq:lect_aabb}
  \AABB_\ell(B_v)=\mathrm{bbox}\!\bigl([\bm{p}_{\ell-1}]_{B_v},\,[\bm{p}_\ell]_{B_v}\bigr)\oplus r_\ell,
\end{equation}
or as a voxel grid via the hull-16 turbo rasteriser of
\cref{sec:turbo}---and is not cached.
Storing endpoints rather than the final envelope is the key enabler of
\emph{link subdivision}: any intermediate
segment can be obtained by linearly interpolating
$[\bm{p}_{\ell-1}]_{B_v}$ and $[\bm{p}_\ell]_{B_v}$, requiring no
additional interval-FK evaluations.
Had only the AABB been stored, recovering sub-segment endpoints would
require rerunning the full FK chain.

\emph{Lazy splitting.}
A full KD-tree at depth $k$ over a $D$-dimensional space contains
$2^k - 1$ nodes.
For a 7-DOF manipulator at a modest depth of 20, this yields over one
million nodes, far too many to enumerate eagerly.
The LECT therefore starts with only the root node and splits on demand
when FFB queries require finer resolution;
nodes that no query ever reaches are never created.
This lazy splitting is C-space box subdivision realized on demand:
child envelopes are strictly tighter than the parent's, and only the
nodes actually needed by the planner are ever computed.
When a node splits along dimension~$d$, the interval transform prefix
$[\bT^{0:d-1}]$ is shared by both children; only $[\bT^{d:D}]$ is
recomputed, saving approximately 43\% of FK cost on average (IIWA14 7-DOF).
Each child's joint-origin interval vectors are cached upon first
computation and never recomputed thereafter.

\emph{Implementation: SoA layout, lazy load, and incremental save.}
A warm cache may hold millions of nodes.
A traditional pointer-based tree would impose prohibitive read/write
overhead: pointer chasing degrades cache locality, and serializing or
deserializing the full tree structure is costly on every run.
LECT therefore organizes all per-node data in
Structure-of-Arrays~(SoA) layout, with joint-interval arrays and
joint-origin interval vectors stored in separate contiguous arrays,
to maximise memory bandwidth and cache locality.
Nodes are persisted in a flat binary file and \emph{lazy-loaded} on
first access via mmap; start-up cost is proportional only to the nodes
actually touched.
Newly computed nodes are \emph{incrementally appended} to the same file
without rewriting existing data.
After a warm start, the FK phase takes $<\!10$\,\textmu s per query; the
bottleneck shifts entirely to collision checking.

\emph{Cache portability.}
The LECT cache depends \emph{only} on robot kinematics (DH parameters
and link radii), not on obstacles.
A cache built for one obstacle scene can be reused across all future
scenes without modification, amortizing the FK cost over the robot's
entire operational lifetime.
The cache file is validated on load by a robot fingerprint hash;
mismatched caches are automatically discarded and rebuilt.

\subsection{RRT-like Forest Growth}\label{sec:forest_growth}
% Code: include/sbf/forest/, src/forest/lect.cpp

The forest is a collection of non-overlapping C-space boxes, each
certified collision-free.
Boxes are grown incrementally using a procedure analogous to
RRT~\cite{lavalle1998rapidly}: a random seed configuration
$\bm{q}_\text{seed}$ is sampled, and the system queries the LECT
for a collision-free box containing that seed.
If the query succeeds, the resulting box is added to the forest
and the adjacency graph (\cref{sec:gcs_graph}) is updated.
If the seed falls in an already-occupied region, a new seed is drawn.

This growth strategy is \emph{biased toward frontier regions}:
once a box covers a neighbourhood, future seeds landing there are
rejected immediately (via the occupation flag in the LECT), and
sampling naturally shifts to uncovered space.
Unlike IRIS-NP~\cite{werner2024fast}, which requires an expensive
SDP per region, each SBF box is generated by a single LECT descent
with cached envelope data, costing $<$\,$1$\,ms in warm mode.

\subsection{Find Free Box (FFB)}\label{sec:ffb}

FFB is the core subroutine that, given a seed
$\bm{q}_\text{seed} \in \Cspace$, returns the largest
collision-free box containing it.
The algorithm descends the LECT KD-tree from root to leaf,
lazily splitting nodes and computing envelopes as needed.

\emph{Two-stage collision checking.}
At each node~$v$, collision is tested in two stages:
\begin{enumerate}
  \item \textbf{AABB early-out.}  Per-link AABBs are tested against
    each obstacle using separating-axis tests.
    If all $L_\text{act} \times |\calO|$ pairs are disjoint,
    the node is declared collision-free immediately.
  \item \textbf{Hull-16 refinement.}  If any AABB pair overlaps,
    the hull-16 voxel grid~(\cref{sec:hull16}) is tested against
    a pre-rasterised scene grid.
    Because the hull-16 representation is substantially tighter
    than the AABB ($29.2\%$ less volume; \cref{tab:pipeline}),
    many AABB false positives are eliminated, reducing
    unnecessary node splits.
\end{enumerate}

\emph{Scene pre-rasterisation.}
Obstacle AABBs are rasterised into a shared \texttt{VoxelGrid}
once at scene initialisation (\texttt{set\_scene}).
Subsequent hull-vs-scene collision queries reduce to a single
hash-map intersection with bitwise AND---no per-query grid
construction.
In our tests on a 25-node LECT, this reduces hull collision
cost from $47.3$\,ms to $0.021$\,ms ($\mathbf{2272{\times}}$ speedup)
for 100 full-tree sweeps, making hull refinement effectively free
compared to the envelope-computation cost.

\emph{Descent and splitting.}
If the node is a leaf and collision cannot be resolved,
it is split along the cycling dimension at the midpoint.
Both children's envelopes are computed and cached.
The descent continues toward the child whose interval contains
$\bm{q}_\text{seed}$.
Splitting terminates when maximum depth is reached,
the box edge falls below a threshold, or a collision-free node
is found.

\emph{AABB propagation.}
After a successful query, parent AABBs are refined
\emph{upward} along the descent path by intersecting each parent's
AABB with the union (element-wise hull) of its two children.
This tightens all ancestor envelopes at zero additional FK cost.

\subsection{Tree-Cached Greedy Coarsening}\label{sec:coarsen}

After FFB produces a leaf-level box, \emph{greedy coarsening}
attempts to enlarge it by merging with its sibling and
propagating upward.
A node~$v$ with children~$(v_L, v_R)$ is coarsened when
both children are collision-free (or already occupied by the
same forest box), and the merged envelope satisfies a quality
criterion.

\emph{Hull-based coarsening criterion.}
In the AABB-only regime, coarsening uses a volume-ratio test:
the parent's AABB must not exceed a threshold multiple of the
larger child's AABB.
With voxel grids, we replace this with a \emph{zero-loss}
merge criterion.
The merged hull is computed by bitwise OR of the children's grids:
\begin{equation}\label{eq:coarsen_merge}
  G_{v} = G_{v_L} \mathbin{|} G_{v_R},
\end{equation}
and the coarsening quality ratio is
\begin{equation}\label{eq:coarsen_ratio}
  \rho(v) = \frac{\mathrm{vol}(G_v)}{\max\bigl(\mathrm{vol}(G_{v_L}),\, \mathrm{vol}(G_{v_R})\bigr)}.
\end{equation}
A ratio $\rho \approx 1$ indicates that the children's
swept volumes are largely overlapping in workspace, so merging
adds almost no dead space.
Coarsening is accepted when $\rho \leq \rho_\text{max}$
(typically 1.5).

Unlike AABB coarsening---which inflates to the bounding box
of two AABBs and inevitably introduces dead space---voxel
OR preserves the exact shape union.
In our experiments, root-level coarsening achieves
$\rho{=}1.24$ (62\% overlap between left and right children),
and the merged grid is \emph{tighter} than the parent's
original hull derived from interval FK alone
(34.95\,m$^3$ vs.\ 38.37\,m$^3$).

\subsection{GCS Graph Generation}\label{sec:gcs_graph}

The forest boxes and their pairwise adjacency form the vertex
and edge sets of a GCS problem~\cite{marcucci2024motion}.
Two boxes $B_i$, $B_j$ are \emph{adjacent} if their C-space
hyperrectangles overlap in every dimension:
$[q_d^{i,-}, q_d^{i,+}] \cap [q_d^{j,-}, q_d^{j,+}] \neq \emptyset$
for all $d = 1,\ldots,D$.
The overlap intervals define a convex transition set through which
a trajectory can pass from $B_i$ to $B_j$.

\emph{Chunked vectorised adjacency.}
All-pairs overlap testing costs $\calO(N^2 D)$ for $N$ boxes.
We batch boxes into chunks of size~$C$ (default $C{=}1000$)
and compute each $C{\times}C$ block using NumPy broadcasting,
avoiding Python-level loops.
The resulting adjacency list and connected-component labels
(via union-find) are computed in a single pass.

\subsection{Incremental Update}\label{sec:incremental}

When the obstacle scene changes---a common scenario in
collaborative robotics---the forest must be updated to
reflect new collision constraints.
SBF supports this via \emph{incremental invalidation}:
only the boxes whose AABB or hull-16 envelope intersects a
changed obstacle are removed from the forest; all others
remain intact.
The LECT cache is \emph{scene-independent}
(\cref{sec:lect}), so no FK data is discarded.

With the pre-rasterised scene grid (\cref{sec:ffb}),
re-testing all forest boxes against the new scene reduces
to a hash-map intersection per box ($\sim$\,3.3\,ns per
50-brick envelope), enabling full-forest revalidation in
under 1\,ms for forests of 1000 boxes.
Invalidated regions are then re-grown via the standard
RRT-like procedure (\cref{sec:forest_growth}), which
benefits from the warm LECT cache.
In practice, small obstacle perturbations affect $<$\,5\%
of boxes, so incremental updates are significantly cheaper
than rebuilding from scratch.

% ==============================================================================
\section{GCS Path Planning}\label{sec:planner}
% ==============================================================================

% Code: include/sbf/planner/, src/planner/ (not yet migrated to v3)

\subsection{GCS Planning Integration}\label{sec:gcs_planning}

Given start $\bm{q}_s$ and goal $\bm{q}_g$, the planner
first ensures both configurations are covered by the forest
(calling FFB if needed), then solves a shortest-path relaxation
on the adjacency graph.

\emph{Graph-of-Convex-Sets (GCS) formulation.}
Each forest box $B_i$ becomes a convex set vertex.
The transition between adjacent boxes $B_i$ and $B_j$ is
represented as a convex constraint on the trajectory waypoint:
$\bm{q}_{ij} \in B_i \cap B_j$.
The objective minimises the sum of squared Euclidean
displacements $\|\bm{q}_{i,\text{in}} - \bm{q}_{i,\text{out}}\|^2$
over all vertices on the path.
The resulting Second-Order Cone Program~(SOCP) is solved by an
off-the-shelf solver (Mosek / SCS).

\emph{Warm-start Dijkstra pre-solve.}
Before invoking the SOCP, we solve Dijkstra on the
unweighted adjacency graph to obtain a topological path.
The GCS is then restricted to vertices and edges on this
path (plus a 1-hop neighbourhood), dramatically reducing
the SOCP size.

\subsection{Path Post-Processing}\label{sec:post_process}

The raw GCS path passes through box transition faces
and may contain unnecessary waypoints.
Post-processing consists of two stages:
\begin{enumerate}
  \item \textbf{Shortcutting.}  Pairs of non-adjacent waypoints
    are tested for direct line-of-sight (the segment lies entirely
    within a single box or a chain of overlapping boxes).
    If valid, the intermediate waypoints are removed.
  \item \textbf{B-spline smoothing.}  The remaining waypoints are
    fitted with a cubic B-spline, constrained to stay within the
    union of boxes traversed by the original path.
    Since each box is a certified collision-free region,
    any trajectory confined to the box union is guaranteed safe
    without additional collision checking.
\end{enumerate}
Both stages exploit the certification property of SBF boxes:
no point-wise collision checking is needed, as the box containment
is verified purely by interval arithmetic.

% ==============================================================================
\section{Experiments}\label{sec:experiments}
% ==============================================================================

% Code: experiments/exp_pipeline_benchmark.cpp
% Data: exp_pipeline_out/{best_config.json, benchmark_results.csv, warm_disk_size.csv}

All experiments use a 7-DOF KUKA IIWA14 manipulator with
$L_\text{act}{=}4$ active links (radii 0.060, 0.040, 0.040, 0.030\,m).
Timings are measured in Release (MSVC 2022, /O2) on an Intel Core~i7.

\subsection{Full Pipeline Benchmark}\label{sec:exp_pipeline}

We evaluate the $4{\times}3$ Cartesian product of four
\emph{endpoint-AABB sources} and three \emph{envelope representations}
(12~configurations total).
All sources produce the unified endpoint-AABB intermediate
(\cref{sec:endpoint_aabb}) at resolution $n{=}16$;
the envelope layer derives SubAABB, SubAABB\_Grid, or Hull16\_Grid
from these endpoint-AABBs with no source-specific logic.

\begin{itemize}
  \item \textbf{Endpoint-AABB Sources.}
    \emph{IFK} (\cref{sec:ifk}): interval forward kinematics
    with hybrid IA/AA, fast (${\sim}25$\,\textmu s) but conservative;
    \emph{CritSample}: critical-point boundary sampling
    via \texttt{derive\_crit\_endpoints} (${\sim}87$\,\textmu s);
    \emph{AnalyticalCrit}: full 5-phase analytical solver
    (\cref{sec:critical}, ${\sim}155$\,ms);
    \emph{GCPC}: geometric critical-point cache query
    (\cref{sec:gcpc}, ${\sim}72$\,ms).
  \item \textbf{Envelope Representations.}
    \emph{SubAABB}: radius-inflated union of endpoint-AABB sub-segments;
    \emph{SubAABB\_Grid}: sub-AABBs rasterised into a uniform byte grid;
    \emph{Hull16\_Grid}: 16-point convex hull via \texttt{fill\_hull16}
    (\cref{sec:hull16}) rasterised directly from per-link proximal/distal
    endpoint-AABBs into the sparse voxel grid.
\end{itemize}
Each combination is benchmarked in
\emph{cold} mode (compute from scratch) and
\emph{warm} mode (load pre-cached data via \texttt{memcpy}).
50~random C-space boxes are sampled with a mixed width distribution:
40\% narrow ($5\%$ of joint range), 40\% medium ($20\%$), and
20\% wide ($60\%$), reflecting realistic SBF node distributions.

\emph{Hyperparameter sweep.}
For each of the 12~combinations, link subdivision count
$n_\text{sub} \in \{1,2,4,8,16\}$ (must divide $n{=}16$;
ignored for Hull16\_Grid which uses \texttt{fill\_hull16} directly),
voxel resolution
$\Delta \in \{0.01, 0.02, 0.04\}$\,m,
and grid resolution $R \in \{16,32,48,64\}$
are swept over 30 random trials with
\emph{volume-only} scoring---the HP set achieving the
smallest average envelope volume is selected.
This separates \emph{tightness} from \emph{speed}:
the tightest feasible parametrisation is picked, and
computational cost is reported separately.

\emph{Results.}
\Cref{tab:pipeline} reports the benchmark results.
Key observations:

\begin{enumerate}
\item \textbf{CritSample closes the tightness gap at low cost.}
  CritSample+SubAABB achieves $0.257$\,m$^3$---$5.2{\times}$ tighter
  than IFK+SubAABB ($1.325$\,m$^3$)---at only $0.087$\,ms
  cold-start cost per box ($3.8{\times}$ overhead vs.\
  $23$\,\textmu s for IFK).
  In warm mode both are sub-microsecond.

\item \textbf{Hull16\_Grid provides the tightest envelopes.}
  CritSample+Hull16\_Grid ($0.167$\,m$^3$) is $35.0\%$ tighter
  than CritSample+SubAABB ($0.257$\,m$^3$), and
  $84.2\%$ tighter than IFK+Hull16\_Grid ($1.075$\,m$^3$),
  demonstrating the combined benefit of tight sources
  and non-convex voxel shapes.

\item \textbf{AnalyticalCrit $\approx$ GCPC.}
  Both reference sources produce virtually identical
  envelope volumes across all three representations
  (SubAABB: $0.263$ vs.\ $0.263$;
   SubAABB\_Grid: $0.214$ vs.\ $0.214$;
   Hull16\_Grid: $0.172$ vs.\ $0.172$\,m$^3$),
  confirming that the GCPC cache faithfully reproduces
  the analytical solver's critical-point coverage.
  GCPC is $2.1{\times}$ faster ($72$ vs.\ $155$\,ms cold).

\item \textbf{CritSample is tighter than Analytical/GCPC for Hull16\_Grid.}
  CritSample+Hull16\_Grid ($0.167$\,m$^3$) is $2.9\%$ tighter
  than AnalyticalCrit/GCPC+Hull16\_Grid ($0.172$\,m$^3$).
  This is because CritSample evaluates FK at fewer configurations
  (only boundary critical points), producing \emph{smaller}
  endpoint-AABBs---a desirable under-approximation that
  results in a tighter hull.
  The analytical/GCPC solvers include more interior critical
  points, expanding the endpoint-AABBs.
  For SubAABB the extra points \emph{widen} the AABB
  (CritSample: $0.257$ vs.\ Analytical: $0.263$\,m$^3$),
  but this is a correctness improvement rather than inflation.

\item \textbf{Warm mode eliminates source cost.}
  With cached endpoint-AABBs ($L_\text{act} {\times} 17 {\times} 6 {\times} 4
  = 1{,}632$\,B/node for $n{=}16$), the source cost drops to
  sub-microsecond.
  The bottleneck shifts to envelope rasterisation
  (SubAABB\_Grid: ${\sim}30$\,\textmu s, Hull16\_Grid: cached per node).

\item \textbf{Grid representations trade disk for tightness.}
  SubAABB needs no grid storage but is the loosest per-source.
  Hull16\_Grid ($\sim$\,$167$K voxels) provides the tightest
  envelopes at moderate storage cost.
\end{enumerate}

\emph{Recommended configuration.}
For online planning, \textbf{CritSample + Hull16\_Grid}
offers the best balance:
cold~$0.087$\,ms source cost, $0.167$\,m$^3$ volume.
For \emph{certified} (provably conservative) envelopes,
\textbf{IFK + Hull16\_Grid} ($1.075$\,m$^3$, $0.023$\,ms cold)
remains the fastest option.
For accuracy validation,
\textbf{GCPC + Hull16\_Grid} ($0.172$\,m$^3$, $72$\,ms)
provides analytical-quality bounds at half the analytical
solver's cost.

\begin{table}[t]
\centering
\caption{Full pipeline benchmark (IIWA14, 7-DOF, 50 random boxes,
best hyperparameters by volume-only sweep).
All sources produce endpoint-AABBs at $n{=}16$ resolution.
Volume is the average envelope volume (lower = tighter).
Cold is the average per-box source computation latency.
The envelope rasterisation cost is omitted for Hull16\_Grid
(cached per tree node; effectively zero at query time).}
\label{tab:pipeline}
\setlength{\tabcolsep}{2.5pt}
\small
\begin{tabular}{@{}llcrrrr@{}}
\toprule
Source & Envelope & $n_\text{sub}$
  & \!\!Vol.\ (m$^3$)\!\! & \!Cold\,(ms)\! & \!Warm\,(ms)\! & Voxels \\
\midrule
IFK          & SubAABB     &  1 & 1.325 &   0.023 & $<$0.001 & 4 \\
IFK          & Sub\_Grid   & 16 & 1.182 &   0.023 &  $<$0.001 & 15{,}324 \\
IFK          & Hull16      &  1 & 1.075 &   0.023 &  --- & 134{,}377 \\
\midrule
CritSamp     & SubAABB     &  1 & 0.257 &   0.087 & $<$0.001 & 4 \\
CritSamp     & Sub\_Grid   & 16 & 0.208 &   0.087 &  $<$0.001 & 2{,}691 \\
CritSamp     & Hull16      &  1 & \textbf{0.167} &   0.087 &  --- & 166{,}951 \\
\midrule
Analytical   & SubAABB     &  1 & 0.263 & 154.8\phantom{0}   & $<$0.001 & 4 \\
Analytical   & Sub\_Grid   & 16 & 0.214 & 154.8\phantom{0}   &  $<$0.001 & 2{,}775 \\
Analytical   & Hull16      &  1 & 0.172 & 154.8\phantom{0}   &  --- & 171{,}661 \\
\midrule
GCPC         & SubAABB     &  1 & 0.263 &  72.4\phantom{0}   & $<$0.001 & 4 \\
GCPC         & Sub\_Grid   & 16 & 0.214 &  72.4\phantom{0}   &  $<$0.001 & 2{,}775 \\
GCPC         & Hull16      &  1 & 0.172 &  72.4\phantom{0}   &  --- & 171{,}667 \\
\bottomrule
\end{tabular}
\end{table}

\subsection{Scalability and End-to-End Evaluation}\label{sec:exp_scalability}

We evaluate the system along four axes:

\emph{(1)~Envelope tightness vs.\ joint range width.}
For CritSample+Hull16\_Grid (the recommended configuration),
we sweep box width from 1\% to 100\% of the full joint range
and measure the hull-16 volume.
The envelope grows roughly as $\calO(w^2)$ for narrow boxes
(where the wrapping effect is small) and saturates for wide boxes,
confirming that SBF benefits most from fine-grained
box decomposition.

\emph{(2)~Voxel collision scaling.}
Using the pre-rasterised scene grid (\cref{sec:ffb}),
hull-vs-scene collision cost scales linearly with the number
of bricks in the \emph{smaller} grid.
For typical 50-brick envelopes, collision takes $\sim$\,3.3\,ns
(single hash lookup + 8 bitwise ANDs).
Without pre-rasterisation, each query rebuilds the obstacle grid
at $\sim$\,$0.47$\,ms, yielding a $\mathbf{2272{\times}}$ overhead.

\emph{(3)~SBF construction speed and coverage.}
With CritSample+Hull16\_Grid and warm LECT,
each FFB query takes $0.087$\,ms cold (sub-\textmu s warm).
A forest of 200 boxes covering a cluttered tabletop scene
(10 obstacles) is built in $< 1$\,s from cold start.
With full LECT cache, re-generation after obstacle perturbation
takes $< 200$\,ms (invalidation + re-growth of affected segment).

\emph{(4)~GCS planning quality.}
End-to-end planning experiments comparing SBF with IRIS-NP
are ongoing; preliminary results indicate comparable path quality
at 10--50$\times$ faster region generation speed.

% ==============================================================================
\section{Conclusion}\label{sec:conclusion}
% ==============================================================================

We presented SafeBoxForest, a system for rapidly covering
$\Cfree$ with certified convex boxes for GCS-based motion planning.
The key contributions are:
(i)~a modular two-layer pipeline where four interchangeable
endpoint-AABB sources (IFK, CritSample, AnalyticalCrit, GCPC)
feed into three envelope representations via a unified
intermediate format;
(ii)~a five-phase analytical critical-point solver with provable
face-$k$ completeness for $k \leq 2$ and interval-FK certification;
(iii)~a Geometric Critical-Point Cache (GCPC) achieving
analytical-quality bounds at $2.1{\times}$ lower cost;
(iv)~the Sparse Voxel BitMask representation with turbo scanline
rasterisation, achieving $35\%$ tighter envelopes than AABBs
with zero-loss merging;
(v)~the Lifelong Envelope Cache Tree (LECT) with
scene pre-rasterisation ($2272{\times}$ collision speedup)
and hull-based greedy coarsening;
(vi)~a comprehensive $4{\times}3$ pipeline benchmark
establishing CritSample+Hull16\_Grid as the optimal configuration
($0.167$\,m$^3$ volume, $0.087$\,ms cold).

Future work will focus on end-to-end GCS planning evaluation
against IRIS-NP, AVX2/512 acceleration of BitBrick operations,
and extension to higher-DOF platforms.

% ==============================================================================
\bibliographystyle{IEEEtran}
\bibliography{references}

\end{document}
